\section{Задание}

\begin{enumerate}[label=\textbf{Task \arabic*:}, leftmargin=*]

\item \textit{[4 балла]}

Выберите датасет и задачу для него (регрессия) --- с Kaggle или \href{https://archive.ics.uci.edu/datasets}{UCI Machine Learning Repository}.

Рекомендации:
\begin{itemize}
    \item количество признаков --- не менее 10
    \item количество объектов --- не менее 100
    \item возраст датасета --- не старее 2010 года
\end{itemize}

Выполните EDA, а именно: придумайте вопросы к данным, которые потребуют использования следующих инструментов (напишите каждый вопрос и код к нему) [\textcolor{DarkOrange}{0,5 балла за покрытие каждого пункта}]:
\begin{itemize}
    \item группировка с агрегацией (groupby) \textit{(минимум два признака задействовано в пункте)};
    \item изменение дискретности по времени (resample) \textit{(можно для этого пункта искусственно добавить время (поминутно $\rightarrow$ в час))};
    \item объединение датафреймов (merge, join или concat) \textit{(можно искусственно разбить датафрейм на два: мальчики/девочки разбили по полу, потом собрали к исходному или другому)};
    \item статическая визуализация \texttt{seaborn};
    \item интерактивный график \texttt{plotly};
    \item визуализация распределений (гистограмма относительных частот, kde, boxplot, violinplot, scatter);
    \item тесты на свойства распределения (например, нормальность);
    \item One Hot Encoding для категориальных признаков (если их нет, то придумайте как некие классы, определяемые по другим числовым признакам).
\end{itemize}

\item \textit{[4 балла]}

\begin{enumerate}[label=\arabic*.]
    \item Выберите один из признаков как таргет (регрессант, предиктант), также выберите некоторе число признаков в качестве регрессоров (предикторов). Обоснуйте, что эти признаки имеет смысл использовать как регрессоры, исходя из предметной области (природы датасета), проведите аналитику --- представьте код и рисунки. [\textcolor{DarkOrange}{1 балл}]

    \textbf{Важно!} Если выполняете исследование данных до разделения на обучающую, которая нужна для моделирования, и отложенную выборки (тестовую), то это потенциал для утечек (data leakage)! Следует сразу отделить часть данных, исследовать их только перед прогоном на уже обученной готовой модели!

    \item Обучите линейную регрессионную модель без регуляризации. Проиллюстрируйте процесс обучения и тестирования модели, запишите выводы о качестве. Метрики: RMSE, MAE, $R^{2}$ [\textcolor{DarkOrange}{1 балл}]

    \item Обучите линейную регрессионную модель с регуляризацией (любой). Для предобработки осуществите удаление выбросов и нормирование --- запишите обоснование для выбора способа поиска выбросов и нормирования. Используйте для этапов Pipeline из sklearn. Проиллюстрируйте процесс обучения и тестирования модели, запишите выводы о качестве. Метрики: RMSE, MAE, $R^{2}$ [\textcolor{DarkOrange}{2 балла}]
\end{enumerate}

\item \textit{[2 балла --- бонус]}

По итогам лекции №1 ответьте на вопросы о месте искусственного интеллекта среди инноваций в крупной компании (пример «Норильского никеля» --- см. видео в группе Telegram):

\begin{enumerate}[label=\arabic*.]
    \item Какие задачи моделирования и управления решаются?
    \item Какие разделы/формальные задачи ИИ задействованы?
    \item Какие физико-химические признаки для них используются?
    \item Укажите (предположите) для нескольких признаков: числовые или категориальные, временная и пространственная дискретность регистрации, сырой признак или агрегат по нескольким, сильно зашумленный, много ли пропусков, особенности закона распределения, потенциал для data leakage.
\end{enumerate}

\item \textit{[3 балла --- бонус]}

Покажите, что из максимизации правдоподобия выборки невязок, распределенных по нормальному закону, следует задача минимизации ошибки MSE, вычисленной по этим невязкам.

Обоснование принесет вам 2 балла.

Если выполните не на листе бумаги, а заполните ячейку (математические символы), то получите ещё 1 балл.

\end{enumerate}